\chapter{Introduction}\label{ch:intro}
%% FUNCTION: Motivates the question and states what the thesis contributes: practice first
%% (round robin works by hand), then theory (EF1 guaranteed, EFX existence open from four agents
%% on, no approximation guarantee for EFX), then the gap, the research question, the
%% contributions and a roadmap. No mathematics in this chapter (chair writing guide).

When heirs divide an estate, flatmates split the furniture of a shared flat or children share a
box of toys, a common procedure is to take turns: everyone picks the item they like best among
those that remain, in a fixed order, round after round. This procedure is called Round Robin. It
needs no computer, no money and no announcement of how much anyone values each item; everyone
only has to make a choice when it is their turn. Variants of it are used in practice, for instance
in a draft at Harvard Business School that allocated courses to students and reversed the order
of choosing in every round \parencite[p.~1065]{budish2011}.

Procedures that do better are available, but they cannot be carried out by hand. A local search
procedure, for instance, finds a fair allocation in the sense studied here in every instance that
has been tested, but it does so by moving goods between bundles thousands of times
\parencite{branzei2025}. For a family at a kitchen table, the realistic alternative to Round Robin
is therefore not the best possible allocation, but dividing the goods somehow. The question is
whether taking turns is good enough.

What counts as fair is a question of theory. The most demanding standard, that nobody prefers
another person's share to their own, often cannot be met when goods are indivisible: with one
painting and two heirs, one of them is left without it. The literature therefore uses weaker
standards. Round Robin always meets a weak one: if someone prefers another person's share, this
preference disappears once a single good, the one they like best, is removed from the other
share \parencite{caragiannis2019,amanatidis2023survey}. This standard is called envy-freeness up
to one good (EF1). A stronger standard requires that the
preference disappears after removing \emph{any} single good from the other share, even the one
they like least; it is called envy-freeness up to any good (EFX). Whether an allocation meeting this stronger standard always exists is one of the
main open questions of the field for four or more people
\parencite[Open Problem~2]{amanatidis2023survey}. Round Robin offers no guarantee for it, not even
an approximate one: in some situations, its allocation misses the stronger standard by an
arbitrarily large factor (Chapter~\ref{ch:background}). A guarantee for the worst case therefore
says nothing useful, and only the typical case can tell whether Round Robin is a good way to reach
the stronger standard.

Existing studies of the typical case of such turn-taking procedures concern other goals. Some
ask which order of turns maximises the expected total or minimum satisfaction
\parencite{bouveret2011,kalinowski2013}. Others prove that, as the number of people grows and their
valuations are independent, Round Robin even reaches the most demanding standard once there are
enough goods, and that a variant which reverses the order of the last round reaches the stronger
standard with high probability \parencite{manurangsi2021}. These results are asymptotic and assume
that people value goods independently of each other. For a few people, for people who largely
agree on what is valuable, and for the question of how far Round Robin misses the stronger standard
when it fails, we found no analysis. \textcite[p.~11]{branzei2025} remarks that Round Robin often
fails it, especially when valuations are correlated, but does not quantify this.

This thesis asks: \emph{how often does Round Robin achieve EFX, and under which conditions is it
therefore a good heuristic for EFX?} We answer the question by simulation. We
vary five conditions one at a time around a common reference case: how much the people agree,
how far apart the values of the goods lie, how many goods each person receives, whether the goods
can be divided evenly, and how many people take part. For each combination, we draw
20{,}000 situations and measure both how often Round Robin succeeds and how far it misses when
it fails.

The thesis makes three contributions. First, it measures how often Round Robin reaches the
stronger standard and how far it misses it, for two to five people and for valuations that range
from independent to identical. Second, it proves three simple statements that delimit these
measurements: Round Robin always succeeds when everyone receives at most two goods and nobody
shares a favourite good; it always fails when everyone ranks the goods identically and the goods
cannot be divided evenly; and when all goods are worth almost the same and can be divided evenly,
it almost always succeeds. Third, it compares Round Robin with dividing the goods at random and
with reversing the last round, and it collects the results in a table that tells people who divide
goods by hand when to rely on Round Robin.

Chapter~\ref{ch:background} defines the fairness standards and explains why Round Robin meets the
weaker but not the stronger one. Chapter~\ref{ch:model} describes the simulation.
Chapter~\ref{ch:results} presents the results one condition at a time, and
Chapter~\ref{ch:random} compares Round Robin with random allocation and with reversing the last
round. Chapter~\ref{ch:discussion} discusses limitations and open questions, and
Chapter~\ref{ch:conclusion} concludes with the decision table.
